\section{Exercise 1. Chooser Option}

\begin{enumerate}
      \item \textbf{Fundamental Theorem of Asset Pricing (FTAP)}

            In a finite discrete-time market, the market is arbitrage-free if and only if there exists an Equivalent Martingale Measure (EMM), denoted by $\mathbb{P}^*$, such that the discounted asset price process $(\tilde{S}_n)_{n \in \mathbb{N}}$ where $\tilde{S}_n = e^{-rn} S_n$ is a martingale under $\mathbb{P}^*$.

            Furthermore, the market is complete if and only if this Equivalent Martingale Measure $\mathbb{P}^*$ is unique. Under this unique EMM, the arbitrage-free price $V_n$ at time $n$ of any attainable $\mathcal{F}_N$-measurable payoff $X$ at maturity $N$ is given by the conditional expectation of its discounted payoff:
            \[ V_n = \mathbb{E}^*\left[ e^{-r(N-n)} X \mid \mathcal{F}_n \right] \]

      \item \textbf{Arbitrage-free price at $T_c$}

            At the decision date $T_c$, the buyer can choose whether the option becomes a Call or a Put. A rational investor will choose the option with the highest value at that time.

            The payoff at time $N$ is given by $(S_N - K)_+ \mathds{1}_{C_{T_c} \ge P_{T_c}} + (K - S_N)_+ \mathds{1}_{C_{T_c} < P_{T_c}}$.

            Using the risk-neutral pricing formula from time $T_c$, the price of the chooser option at $T_c$, let's denote it $V_{T_c}$, is:
            \[ V_{T_c} = \mathbb{E}^*\left[ e^{-r(N-T_c)} \left( (S_N - K)_+ \mathds{1}_{C_{T_c} \ge P_{T_c}} + (K - S_N)_+ \mathds{1}_{C_{T_c} < P_{T_c}} \right) \mid \mathcal{F}_{T_c} \right] \]

            Since the indicator functions $\mathds{1}_{C_{T_c} \ge P_{T_c}}$ and $\mathds{1}_{C_{T_c} < P_{T_c}}$ are $\mathcal{F}_{T_c}$-measurable, we can take them out of the conditional expectation:
            \[ V_{T_c} = \mathds{1}_{C_{T_c} \ge P_{T_c}} \mathbb{E}^*\left[ e^{-r(N-T_c)} (S_N - K)_+ \mid \mathcal{F}_{T_c} \right] + \mathds{1}_{C_{T_c} < P_{T_c}} \mathbb{E}^*\left[ e^{-r(N-T_c)} (K - S_N)_+ \mid \mathcal{F}_{T_c} \right] \]

            Recognizing the pricing formulas for the Call and Put options at time $T_c$:
            \[ C_{T_c} = \mathbb{E}^*\left[ e^{-r(N-T_c)} (S_N - K)_+ \mid \mathcal{F}_{T_c} \right] \]
            \[ P_{T_c} = \mathbb{E}^*\left[ e^{-r(N-T_c)} (K - S_N)_+ \mid \mathcal{F}_{T_c} \right] \]

            We get:
            \[ V_{T_c} = \mathds{1}_{C_{T_c} \ge P_{T_c}} C_{T_c} + \mathds{1}_{C_{T_c} < P_{T_c}} P_{T_c} \]

            This expression is exactly the definition of the maximum between $C_{T_c}$ and $P_{T_c}$. Thus, $V_{T_c} = \max(C_{T_c}, P_{T_c})$.

      \item \textbf{Price at $T_c$ using Put-Call parity}

            The Put-Call parity relation at time $t$ for European options on a non-dividend-paying stock is given by:
            \[ C_t - P_t = S_t - e^{-r(N-t)} K \]

            Rearranging this formula to express the Put price in terms of the Call price at time $T_c$:
            \[ P_{T_c} = C_{T_c} - S_{T_c} + e^{-r(N-T_c)} K \]

            Now, substitute this expression for $P_{T_c}$ into the result from question 2:
            \[ V_{T_c} = \max(C_{T_c}, P_{T_c}) = \max\left(C_{T_c}, C_{T_c} - S_{T_c} + e^{-r(N-T_c)} K\right) \]

            We can factor out $C_{T_c}$ from the maximum function:
            \[ V_{T_c} = C_{T_c} + \max\left(0, e^{-r(N-T_c)} K - S_{T_c}\right) \]

            This is equivalent to:
            \[ V_{T_c} = C_{T_c} + \left(e^{-r(N-T_c)} K - S_{T_c}\right)_+ \]

      \item \textbf{Replication portfolio}

            From the previous question, the value of the "chooser" option at time $T_c$ is $C_{T_c} + (e^{-r(N-T_c)} K - S_{T_c})_+$. This value can be viewed as the payoff of a specific portfolio at time $T_c$.

            Let's break down the two terms:
            \begin{itemize}
                  \item $C_{T_c}$ is the value at time $T_c$ of a standard European Call option with strike $K$ and maturity $N$.
                  \item $(e^{-r(N-T_c)} K - S_{T_c})_+$ is exactly the payoff at time $T_c$ of a standard European Put option with strike $K' = e^{-r(N-T_c)} K$ and maturity $T_c$.
            \end{itemize}

            Therefore, to replicate the payoff of the chooser option at $T_c$, an investor can buy a portfolio at time 0 consisting of:
            \begin{itemize}
                  \item One European Call option with strike $K$ and maturity $N$.
                  \item One European Put option with strike $e^{-r(N-T_c)} K$ and maturity $T_c$.
            \end{itemize}

      \item \textbf{Arbitrage-free price at 0}

            By the principle of no-arbitrage, if two portfolios have the same value at a future date $T_c$ in all states of the world, they must have the same value today (time 0).

            Since the chooser option can be perfectly replicated by the portfolio described in question 4, its price at time 0 is simply the sum of the prices of the components of that replicating portfolio evaluated at time 0.

            Using the given notation:
            \begin{itemize}
                  \item The price at 0 of the Call option with strike $K$ and maturity $N$ is $C_0(K, N)$.
                  \item The price at 0 of the Put option with strike $e^{-r(N-T_c)} K$ and maturity $T_c$ is $P_0(e^{-r(N-T_c)} K, T_c)$.
            \end{itemize}

            Thus, the arbitrage-free price of the chooser option at time 0 is:
            \[ V_0 = C_0(K, N) + P_0\left(e^{-r(N-T_c)} K, T_c\right) \]

\end{enumerate}